\begin{proposition}[Proved: a two-generator candidate has PD torsion]
\label{B:PD-two-generator-rejected}
For every prime $p$, the ring
$R_1=\mathbb F_p[x,y,z]/(xy,xz)$ is not PD torsion free.
More precisely, in the universal PD envelope $D$ of
$(p,xy,xz)\subset\mathbb Z_p[x,y,z]$, with the canonical divided
powers on $p$, the element
\[
 \tau_1=z^{p^2-p}\gamma_{p^2}(xy)\gamma_p(xz)
       -y^{p^2-p}\gamma_p(xy)\gamma_{p^2}(xz)
\]
is nonzero, is killed by $p$, and remains nonzero in the $p$-adic
completion of $D$.
\end{proposition}

\begin{proof}
Put $P=\mathbb Z_p[x,y,z]$, and write $U^{[n]}$, $V^{[n]}$ for
basis divided powers in the PD polynomial algebra $P\langle U,V\rangle$.
The universal presentation of \cite[Lemma 60.2.4, Tag 07HC]{B:StacksPDPresentation}
gives
\begin{equation}\label{B:two-generator-presentation}
 D=P\langle U,V\rangle/
 \bigl(U-xy,V-xz,z^qU^{[q]}-y^qV^{[q]}\ (q\geq1)\bigr).
\end{equation}
We verify that no relations have been omitted. If
$rxy+sxz=r_0\in pP$, reduction modulo $p$ gives
$\overline r\,y+\overline s\,z=0$, since $x$ is a nonzerodivisor.
The syzygies of $y,z$ give
$r=zh+pa$ and $s=-yh+pb$ for some $h,a,b\in P$, and then
$r_0=px(ay+bz)$. The relation to which the cited presentation applies is
\[
 rU+sV-r_0=h(zU-yV)+pa(U-xy)+pb(V-xz).
\]
Every positive divided power of the two $p$-multiples belongs to
the ordinary ideal $(U-xy,V-xz)$: for example it is
$p^na^n(U-xy)^n/n!$, with integral coefficient $p^n/n!$.
The addition formula therefore reduces all the required relations
to $\gamma_q(zU-yV)$, $q\geq1$. The addition formula, or equivalently
multiplication of the divided-power generating series by the inverse
series for $yV$, shows inductively that these generate the same
ordinary ideal as $z^qU^{[q]}-y^qV^{[q]}$. This proves
\eqref{B:two-generator-presentation} with the prescribed base divided powers.

Set $s=p^2$ and $r=p$. Multiplying the relation of index $q=s-r$
in \eqref{B:two-generator-presentation} by $U^{[r]}V^{[r]}$ gives
\[
 \binom{p^2}{p}\tau_1=0.
\]
The binomial coefficient has $p$-adic valuation one, as follows
immediately by applying the factorial valuation formula to
$(p^2)!$, $p!$ and $(p^2-p)!$. It is $p$ times a unit in
$\mathbb Z_p$, so $p\tau_1=0$.

For nonvanishing, give $x,y,z$ their standard multidegrees and set
$\deg U=(1,1,0)$, $\deg V=(1,0,1)$. The component of degree
$(s+r,s,s)$ of the free PD polynomial algebra has the
$\mathbb Z_p$-basis
\[
 e_{m,n}=x^{s+r-m-n}y^{s-m}z^{s-n}U^{[m]}V^{[n]},
 \qquad 0\leq m,n\leq s,\quad m+n\leq s+r.
\]
Its quotient relations are generated by all the homogeneous
multiples of the three families in
\eqref{B:two-generator-presentation}. In this basis they are
\begin{align*}
 &(m+1)e_{m+1,n}-e_{m,n},\\
 &(n+1)e_{m,n+1}-e_{m,n},\\
 &\binom{m+h}{h}e_{m+h,n}
       -\binom{n+h}{h}e_{m,n+h}\qquad(h\geq1),
\end{align*}
whenever the indicated terms occur. In the third family the
necessary and sufficient index bounds are
$m+h\leq s$, $n+h\leq s$, and $m+n+h\leq s+r$.
Modulo $p$, take the coefficient functional $\lambda$ at the
corner $(m,n)=(s,r)$. Incoming linear relations have coefficients
$s$ or $r$, both divisible by $p$. There is no outgoing linear
relation from this corner, because $m+n=s+r$ is maximal.
If the corner is the first term of a scalar syzygy, then
$0<h\leq s-r<s$ and its coefficient is
$\binom{s}{h}\equiv0\pmod p$, since $s=p^2$.
It cannot be the second term, since that would require both
$m=s$ and $m+h\leq s$ with $h>0$.
Thus $\lambda$ annihilates every relation, whereas
\[
 \tau_1=e_{s,r}-e_{r,s},\qquad \lambda(\tau_1)=1.
\]
This proves nonvanishing already in $D/pD$.
By \cite[Lemma 10.96.3, Tag 05GG]{B:StacksCompletion}, reduction
modulo $p$ is unchanged by $p$-adic completion, so the image of
$\tau_1$ there is still nonzero. This is a legitimate polynomial
presentation in the PD torsion-freeness criterion, proving the assertion.

The supplementary executable \texttt{B\_finish\_PD\_two\_generator.py}
checks every homogeneous relation in this component for $p=2,3,5,7$,
with respective basis sizes $22,79,466,1597$; its saved JSON records
$\lambda(\tau_1)=1$ and valuation one of the torsion multiplier in
each case. The preceding proof treats every prime and establishes
that the finite relation lists are complete. This proposition is
only a PD-torsion obstruction for a rejected positive candidate;
no horizontality assertion about $\tau_1$ is used.
\end{proof}
